% =====================================================================
% datos.tex
% Responsabilidad única: describir el contrato de datos, el momento de cada
% columna, el tamaño de las tablas y el origen (simulado) de los datos, con
% la forma de conectar la exportación real. Cifras: macros de 01_datos.
% =====================================================================
\section{Datos}
\label{sec:datos}

Antes de modelar nada necesitábamos dos cosas. Por un lado saber qué columnas
tiene la base y en qué momento del proceso se conoce cada una y por otro
tener datos con los que trabajar mientras llegaba la exportación real. Esta sección cuenta cómo
resolvimos las dos.

\subsection{Contrato de datos y momento de cada columna}

La base la describimos con un contrato (\codigo{proyecto/datos/esquema.py}),
que es el único lugar del código donde aparecen los nombres de las columnas.
El resto del proyecto nunca escribe una lista de nombres a mano sino que se
la pide al contrato así que si mañana cambia un nombre se cambia en un solo
sitio. El contrato declara \DatTablasContrato{} tablas (clientes, muestreos,
muestras y controles de calidad) con \DatColumnasContrato{} columnas, de las
que \DatColumnasObligatorias{} son obligatorias (sin ellas el proyecto no
corre) y \DatColumnasOpcionales{} opcionales (si faltan, los análisis que las
usan se saltan solos). A esas tablas se suma un calendario con los festivos
de Colombia.

Cada columna lleva además un \emph{momento}, que es la etapa del proceso en
la que ese dato ya se conoce: \codigo{maestro} (los atributos del cliente),
\codigo{programacion}, \codigo{toma}, \codigo{llegada}, \codigo{recepcion},
\codigo{siembra}, \codigo{lectura}, \codigo{reensayo} y \codigo{reporte}, en
total \DatMomentos{} momentos en orden (Tabla~\ref{tab:momentos}). De ahí sale
una regla sencilla: Una variable solo puede entrar a un modelo si su momento
es anterior o igual al momento en que se hace la predicción. Por ejemplo el
clasificador de riesgo predice en \codigo{recepcion} y por eso no puede ver
el recuento de la siembra ni la fecha del informe, aunque las dos columnas
estén en la misma fila de la tabla. A esto se le llama evitar una
\emph{fuga de información}, que es cuando el modelo aprende de algo que en la
práctica todavía no se sabría. Las columnas que agrega la limpieza (como
\codigo{dias\_entrega} o \codigo{no\_conforme}) también declaran su momento y
cada tabla de modelado pasa por una auditoría automática de fugas antes de
usarse (Sección~\ref{sec:cribado}).

\begin{table}[!htbp]
\centering
\begin{minipage}[t]{0.40\textwidth}\centering\small
\tablagen{filas_crudas}
\end{minipage}\hfill
\begin{minipage}[t]{0.40\textwidth}\centering\small
\tablagen{columnas_por_momento}
\end{minipage}
\caption{Filas de cada tabla cruda (izquierda) y columnas del contrato que se
conocen en cada momento del proceso (derecha), en el orden en que ocurren: Un
modelo que predice en un momento solo ve las columnas de su fila y de las
anteriores.}
\label{tab:momentos}
\end{table}

En total leímos \DatMuestrasCrudas{} filas crudas de muestras de
\DatClientes{} clientes, \DatVisitas{} visitas de muestreo en planta y
\DatControles{} controles de calidad (positivos y negativos) entre
\DatDesde{} y \DatHasta{}. El origen trae todas las columnas opcionales
(faltan \DatFaltanOpcionales{}), lo cual es de esperar en una simulación. Con
la exportación real puede que falten algunas y para eso la etapa~01 produce
una tabla de cobertura que dice de antemano qué análisis quedarían por fuera.

\subsection{Base real y complemento simulado}
\label{sec:simulacion}

La base del proyecto es la exportación real del laboratorio
(Sección~\ref{sec:datos-reales}). Como no trae los controles de calidad de
los ensayos ni los datos de recepción, que varias técnicas del curso
necesitan, la complementamos con una base simulada de la misma estructura
(\codigo{proyecto/datos/simulacion.py}, semilla \DatSemilla{}) y sobre ella
aplicamos lo de las Secciones~\ref{sec:exploracion} a~\ref{sec:evaluacion}.
La base complementaria nos da además algo que la real no puede dar: La
verdad detrás de cada muestra, con la que podemos auditar lo que encuentran
los modelos. El generador produce
tablas con la misma estructura del contrato y con el tipo de suciedad que
esperamos de una exportación real (faltantes, categorías mal escritas,
duplicados, valores imposibles y recuentos escritos como texto). El conjunto
de datos simulado es entonces una herramienta del trabajo y no su aporte.

Quisimos que el generador no favoreciera a ningún modelo por construcción y
por eso la señal mezcla efectos lineales, umbrales,
interacciones y episodios que solo pasan en ciertos meses. Además hay
columnas que no aportan nada y otras redundantes. Las variables que no se
observan en la realidad (si la muestra estaba de verdad contaminada, la
higiene y el segmento de cada cliente o el error de cada analista) no se
escriben en las tablas sino aparte, en una ``verdad oculta'' que solo
consultan las celdas de auditoría marcadas como tales. Ningún módulo por
fuera de \codigo{proyecto/datos/} importa el generador, de modo que el resto
del código no sabe si está trabajando con datos simulados o reales.

Para conectar la exportación real bastaría con copiar los archivos en
\codigo{datos/reales/}, completar \codigo{configuracion/mapeo\_columnas.json}
(el nombre real de cada columna y la traducción de las etiquetas) y
reconstruir los notebooks con el origen \codigo{real}. Las celdas de
auditoría de la simulación se saltan solas y la huella de las tablas crudas
(\codigo{\DatHuella}) cambiaría, lo que deja constancia de con qué datos se
produjo cada cifra.
